\section{Edit Distance}
\label{sec:dp-edit-distance}

\problemstatement{Given two strings $s_1$ (length $n$) and $s_2$ (length
$m$), find the minimum number of operations -- \textbf{insert},
\textbf{delete}, or \textbf{replace} a single character -- needed to
transform $s_1$ into $s_2$.}

\begin{patternbox}
This is Section~\ref{sec:dp-lcs}'s match/mismatch template, extended:
the \textbf{match} case is unchanged (no operation needed, advance both
pointers for free), but the \textbf{mismatch} case is no longer a simple
``discard one side'' -- it now branches into \textbf{three} possible
operations, each costing $1$ and corresponding to a different way of
advancing the two pointers.
\end{patternbox}

\subsection*{Deriving the recurrence}

\[
  f(i,j) =
  \begin{cases}
    f(i{-}1,j{-}1) & \text{if } s_1[i{-}1]=s_2[j{-}1], \\
    1 + \min\big(f(i{-}1,j{-}1),\ f(i{-}1,j),\ f(i,j{-}1)\big) & \text{otherwise.}
  \end{cases}
\]

On a mismatch, three operations are available, each mapping to a
specific pointer movement: \textbf{replace} $s_1[i-1]$ with
$s_2[j-1]$ (both characters consumed, advance both pointers --
$f(i{-}1,j{-}1)$); \textbf{delete} $s_1[i-1]$ (advance only $s_1$'s
pointer, leaving $s_2[j-1]$ still to be matched -- $f(i{-}1,j)$); or
\textbf{insert} a copy of $s_2[j-1]$ into $s_1$ (this newly-inserted
character matches $s_2[j-1]$ for free, so only $s_2$'s pointer advances
-- $f(i,j{-}1)$). We pay $1$ for whichever operation is used, plus the
cost of solving the resulting smaller sub-problem, and take the cheapest
option.

Base cases: $f(0,j)=j$ (transform an empty $s_1$ into $s_2$'s first $j$
characters by $j$ insertions), $f(i,0)=i$ (transform $s_1$'s first $i$
characters into an empty string by $i$ deletions).

\subsection*{Approach 1: Recursion}

\begin{lstlisting}
#include <bits/stdc++.h>
using namespace std;

int editDistanceRecursive(int i, int j, string& s1, string& s2) {
    if (i == 0) return j;
    if (j == 0) return i;

    if (s1[i - 1] == s2[j - 1]) {
        return editDistanceRecursive(i - 1, j - 1, s1, s2);
    }
    int replace = editDistanceRecursive(i - 1, j - 1, s1, s2);
    int del     = editDistanceRecursive(i - 1, j, s1, s2);
    int insert  = editDistanceRecursive(i, j - 1, s1, s2);
    return 1 + min({replace, del, insert});
}

int main() {
    string s1 = "horse", s2 = "ros";
    cout << editDistanceRecursive(s1.size(), s2.size(), s1, s2) << endl; // 3
    return 0;
}
\end{lstlisting}

\begin{complexitybox}[Approach 1 Complexity]
\textbf{Time.} $O(3^{n+m})$ in the worst case -- three branches per
mismatch.
\textbf{Space.} $O(n+m)$ recursion stack.
\end{complexitybox}

\subsection*{Approach 2: Memoization}

\begin{lstlisting}
#include <bits/stdc++.h>
using namespace std;

int editDistanceMemo(int i, int j, string& s1, string& s2,
                      vector<vector<int>>& memo) {
    if (i == 0) return j;
    if (j == 0) return i;
    if (memo[i][j] != -1) return memo[i][j];

    if (s1[i - 1] == s2[j - 1]) {
        return memo[i][j] = editDistanceMemo(i - 1, j - 1, s1, s2, memo);
    }
    int replace = editDistanceMemo(i - 1, j - 1, s1, s2, memo);
    int del     = editDistanceMemo(i - 1, j, s1, s2, memo);
    int insert  = editDistanceMemo(i, j - 1, s1, s2, memo);
    return memo[i][j] = 1 + min({replace, del, insert});
}

int main() {
    string s1 = "horse", s2 = "ros";
    int n = s1.size(), m = s2.size();
    vector<vector<int>> memo(n + 1, vector<int>(m + 1, -1));
    cout << editDistanceMemo(n, m, s1, s2, memo) << endl; // 3
    return 0;
}
\end{lstlisting}

\begin{complexitybox}[Approach 2 Complexity]
\textbf{Time.} $O(n \cdot m)$.
\textbf{Space.} $O(n \cdot m)$ memo $+$ $O(n+m)$ recursion stack.
\end{complexitybox}

\subsection*{Approach 3: Tabulation}

\begin{lstlisting}
#include <bits/stdc++.h>
using namespace std;

int editDistanceTabulation(string& s1, string& s2) {
    int n = s1.size(), m = s2.size();
    vector<vector<int>> dp(n + 1, vector<int>(m + 1, 0));

    for (int j = 0; j <= m; j++) dp[0][j] = j;
    for (int i = 0; i <= n; i++) dp[i][0] = i;

    for (int i = 1; i <= n; i++) {
        for (int j = 1; j <= m; j++) {
            if (s1[i - 1] == s2[j - 1]) {
                dp[i][j] = dp[i - 1][j - 1];
            } else {
                dp[i][j] = 1 + min({dp[i - 1][j - 1], dp[i - 1][j], dp[i][j - 1]});
            }
        }
    }
    return dp[n][m];
}

int main() {
    string s1 = "horse", s2 = "ros";
    cout << editDistanceTabulation(s1, s2) << endl; // 3
    return 0;
}
\end{lstlisting}

\subsection*{Dry run of the tabulation table}

\dryrun{Take $s_1=\texttt{"horse"}$, $s_2=\texttt{"ros"}$.}

\begin{center}
\begin{tabular}{c|c|c|c|c}
 & \texttt{""} & \texttt{r} & \texttt{o} & \texttt{s} \\ \hline
\texttt{""} & 0 & 1 & 2 & 3 \\ \hline
\texttt{h} & 1 & 1 & 2 & 3 \\ \hline
\texttt{o} & 2 & 2 & 1 & 2 \\ \hline
\texttt{r} & 3 & 2 & 2 & 2 \\ \hline
\texttt{s} & 4 & 3 & 3 & 2 \\ \hline
\texttt{e} & 5 & 4 & 4 & 3
\end{tabular}
\end{center}

The answer is $\textit{dp}[5][3]=3$: one optimal sequence is
\texttt{"horse"} $\to$ \texttt{"rorse"} (replace \texttt{h}$\to$\texttt{r})
$\to$ \texttt{"rose"} (delete the second \texttt{r}) $\to$
\texttt{"ros"} (delete \texttt{e}): $3$ operations.

\begin{complexitybox}[Approach 3 Complexity]
\textbf{Time.} $O(n \cdot m)$.
\textbf{Space.} $O(n \cdot m)$ table, no recursion stack.
\end{complexitybox}

\subsection*{Approach 4: Space Optimization}

\begin{lstlisting}
#include <bits/stdc++.h>
using namespace std;

int editDistanceSpaceOptimized(string& s1, string& s2) {
    int n = s1.size(), m = s2.size();
    vector<int> prev(m + 1), current(m + 1);
    for (int j = 0; j <= m; j++) prev[j] = j;

    for (int i = 1; i <= n; i++) {
        current[0] = i;
        for (int j = 1; j <= m; j++) {
            if (s1[i - 1] == s2[j - 1]) {
                current[j] = prev[j - 1];
            } else {
                current[j] = 1 + min({prev[j - 1], prev[j], current[j - 1]});
            }
        }
        prev = current;
    }
    return prev[m];
}

int main() {
    string s1 = "horse", s2 = "ros";
    cout << editDistanceSpaceOptimized(s1, s2) << endl; // 3
    return 0;
}
\end{lstlisting}

\begin{complexitybox}[Approach 4 Complexity]
\textbf{Time.} $O(n \cdot m)$.
\textbf{Space.} $O(m)$ for two rolling rows.
\end{complexitybox}

\begin{takeawaybox}
Edit Distance generalizes the LCS template's mismatch case from ``pick
the better of two discards'' to ``pick the cheapest of three specific
operations,'' each mapping to a distinct, intuitive pointer movement.
Once this three-way branch is understood, Delete Operations
(Section~\ref{sec:dp-delete-operations}) is recoverable as the special
case where only the delete operation is permitted.
\end{takeawaybox}
